% Options for packages loaded elsewhere
\PassOptionsToPackage{unicode}{hyperref}
\PassOptionsToPackage{hyphens}{url}
\documentclass[
  ignorenonframetext,
]{beamer}
\newif\ifbibliography
\usepackage{pgfpages}
\setbeamertemplate{caption}[numbered]
\setbeamertemplate{caption label separator}{: }
\setbeamercolor{caption name}{fg=normal text.fg}
\beamertemplatenavigationsymbolsempty
% remove section numbering
\setbeamertemplate{part page}{
  \centering
  \begin{beamercolorbox}[sep=16pt,center]{part title}
    \usebeamerfont{part title}\insertpart\par
  \end{beamercolorbox}
}
\setbeamertemplate{section page}{
  \centering
  \begin{beamercolorbox}[sep=12pt,center]{section title}
    \usebeamerfont{section title}\insertsection\par
  \end{beamercolorbox}
}
\setbeamertemplate{subsection page}{
  \centering
  \begin{beamercolorbox}[sep=8pt,center]{subsection title}
    \usebeamerfont{subsection title}\insertsubsection\par
  \end{beamercolorbox}
}
% Prevent slide breaks in the middle of a paragraph
\widowpenalties 1 10000
\raggedbottom
\AtBeginPart{
  \frame{\partpage}
}
\AtBeginSection{
  \ifbibliography
  \else
    \frame{\sectionpage}
  \fi
}
\AtBeginSubsection{
  \frame{\subsectionpage}
}
\usepackage{iftex}
\ifPDFTeX
  \usepackage[T1]{fontenc}
  \usepackage[utf8]{inputenc}
  \usepackage{textcomp} % provide euro and other symbols
\else % if luatex or xetex
  \usepackage{unicode-math} % this also loads fontspec
  \defaultfontfeatures{Scale=MatchLowercase}
  \defaultfontfeatures[\rmfamily]{Ligatures=TeX,Scale=1}
\fi
\usepackage{lmodern}
\ifPDFTeX\else
  % xetex/luatex font selection
\fi
% Use upquote if available, for straight quotes in verbatim environments
\IfFileExists{upquote.sty}{\usepackage{upquote}}{}
\IfFileExists{microtype.sty}{% use microtype if available
  \usepackage[]{microtype}
  \UseMicrotypeSet[protrusion]{basicmath} % disable protrusion for tt fonts
}{}
\makeatletter
\@ifundefined{KOMAClassName}{% if non-KOMA class
  \IfFileExists{parskip.sty}{%
    \usepackage{parskip}
  }{% else
    \setlength{\parindent}{0pt}
    \setlength{\parskip}{6pt plus 2pt minus 1pt}}
}{% if KOMA class
  \KOMAoptions{parskip=half}}
\makeatother
\setlength{\emergencystretch}{3em} % prevent overfull lines
\providecommand{\tightlist}{%
  \setlength{\itemsep}{0pt}\setlength{\parskip}{0pt}}
\usepackage{bookmark}
\IfFileExists{xurl.sty}{\usepackage{xurl}}{} % add URL line breaks if available
\urlstyle{same}
\hypersetup{
  hidelinks,
  pdfcreator={LaTeX via pandoc}}

\author{\texorpdfstring{}{}}
\date{}

\begin{document}

\section{Abstract}\label{abstract}

\begin{frame}[allowframebreaks]{Abstract}
The Cognitive Integrity Framework (CIF) \cite{friedman2026cogsec1}
establishes formal guarantees for the safe operation of multiagent AI
systems through five canonical defense mechanisms: Cognitive Firewalls,
Belief Sandboxing, Behavioral Invariants, Drift Detection, and Byzantine
Consensus. Paper 2 of this series \cite{friedman2026cogsec2}
demonstrated 94--100\% detection at the parametric ceiling across 950
attack scenarios and four production multiagent architectures, and Paper
3 \cite{friedman2026cogsec3} translated these results into accessible
deployment guidance, incident-response playbooks, and operator risk
frameworks. However, the practical applicability of these mechanisms
across diverse operational domains has not been systematically
evaluated. This paper---Part 4 of the \emph{Cognitive Security for
Multiagent Operators} series---addresses this gap.

We present the \textbf{CIF-AD-OODA integration model}, combining CIF's
defense mechanisms with Axiomatic Design (AD) theory
\cite{suh2001axiomatic} and Boyd's OODA (Observe-Orient-Decide-Act) Loop
\cite{boyd1987patterns} to analyze \textbf{Goal Hijacking}: an advanced
form of indirect prompt injection that constitutes a \emph{teleological
attack} on autonomous agency. Unlike conventional exploits targeting
data confidentiality, Goal Hijacking corrupts the OODA Loop's
\textbf{Orientation} phase, overriding an agent's Functional
Requirements (FRs) while preserving the appearance of rational behavior.
We model these attacks as ``fast OODA transients'' that introduce
off-diagonal coupling in the agent's Design Matrix, violating Suh's
Independence Axiom.

Using a standardized five-step domain analysis template, we conduct a
rigorous analysis across ten critical domains: (1) Rare Earth mining,
(2) Nation-state alliances, (3) Cyber-security, (4) Drone warfare, (5)
Supply chains, (6) Biowarfare, (7) Food security, (8) Trade wars, (9)
Infrastructure, and (10) Information ecosystems. Cross-domain synthesis
reveals three universal attack patterns---FR Polarity Inversion,
Constraint Relaxation, and Context Boundary Violation---and confirms
that all five CIF mechanisms provide adequate coverage across the domain
portfolio. Three novel defense extensions emerge: \emph{verification
channel separation} (Biowarfare), \emph{active perturbation probing}
(Trade Wars), and \emph{physics-informed invariants} (Infrastructure).
Retrospective analysis of six documented AI agent security incidents
(2024--2025)---including production database destruction, remote code
execution via prompt injection, and \$3.2M procurement fraud---validates
the attack pattern taxonomy and confirms CIF defense applicability.
These findings are corroborated by the OWASP Top 10 for Agentic
Applications (2025) \cite{owasp2025agentic}, which designates Agent Goal
Hijack as the primary risk for deployed agent systems. We demonstrate
that without CIF-mediated \textbf{Drift Detection} and
\textbf{Behavioral Invariants} stabilizing the Orientation phase,
high-efficiency agents are liable to optimize for adversarial objectives
under the guise of compliance.

\textbf{Paper Series.} This is Part 4 of the four-part \emph{Cognitive
Security for Multiagent Operators} series: Part 1 (DOI:
10.5281/zenodo.18364119) develops the formal foundations (trust
calculus, defense composition algebra, adversary taxonomy,
information-theoretic stealth--impact bounds, and model-checked
invariants); Part 2 (DOI: 10.5281/zenodo.18364128) provides
computational validation across production multiagent architectures;
Part 3 (DOI: 10.5281/zenodo.18364130) presents a qualitative review and
practitioner's deployment guide; Part 4 (this paper) applies the
framework across ten critical operational domains through the integrated
CIF-AD-OODA model. Readers seeking derivations should consult Part 1;
readers seeking empirical measurements, Part 2; readers seeking
deployment guidance, Part 3.
\end{frame}

\end{document}
